\documentclass[11pt]{article}

\usepackage[utf8]{inputenc}
\usepackage[T1]{fontenc}
\usepackage{geometry}
\usepackage{amsmath}
\usepackage{amssymb}
\usepackage{booktabs}
\usepackage{hyperref}
\usepackage{xurl}
\usepackage{xcolor}
\usepackage{listings}
\usepackage{enumitem}
\usepackage{parskip}
\usepackage{graphicx}
\graphicspath{{../images/}}

\geometry{margin=1in}

\lstset{
  language=Java,
  basicstyle=\ttfamily\small,
  keywordstyle=\color{blue},
  commentstyle=\color{gray},
  stringstyle=\color{teal},
  showstringspaces=false,
  columns=fullflexible,
  frame=single,
  framerule=0.2pt,
  breaklines=true
}

\setlist[itemize]{topsep=0.3em,itemsep=0.2em}
\setlist[enumerate]{topsep=0.3em,itemsep=0.2em}

% Simple per-section source line.
\newcommand{\notesources}[1]{\par\noindent{\small\color{gray}\textit{Sources:} \sloppy #1\par}}

% Unit-wise source strings for this subject (used by slides/notes).
% Path is relative to the lecture `latex/` folder (the compilation CWD).
% OOPS with Java (CSEG1044) — unit-wise sources
%
% These are short, student-facing reference strings intended to be used with:
%   - \slidesources{...}  (in beamer slides)
%   - \notesources{...}   (in lecture notes)
%
% Style inspiration: other subjects in My-Subjects/ use semicolon-separated sources
% and include an "accessed" date for web documentation.

\newcommand{\OOPSUnitISources}{Schildt, \textit{Java: A Beginner's Guide} (10e), McGraw-Hill (Java fundamentals + OOP basics).; Oracle Java Tutorials: Getting Started + Language Basics + Classes and Objects (accessed \today).; Java SE API docs (e.g., \texttt{String}, \texttt{Scanner}) (accessed \today).}

\newcommand{\OOPSUnitIISources}{Schildt, \textit{Java: The Complete Reference} (12e), McGraw-Hill (classes, inheritance, packages, interfaces).; Bloch, \textit{Effective Java} (3e), Addison-Wesley, 2018 (design choices; interfaces vs abstract classes).; Oracle Java Tutorials: Classes and Objects + Interfaces + Packages (accessed \today).; Java SE API docs (e.g., \texttt{Comparable}, \texttt{Runnable}) (accessed \today).}

\newcommand{\OOPSUnitIIISources}{Schildt, \textit{Java: The Complete Reference} (12e) (nested/inner classes, exceptions, threads, I/O).; Oracle Java Tutorials: Nested Classes + Exceptions + Concurrency + Basic I/O (accessed \today).; Java SE API docs (e.g., \texttt{Thread}, \texttt{AutoCloseable}) (accessed \today).}

\newcommand{\OOPSUnitIVSources}{Schildt, \textit{Java: The Complete Reference} (12e) (generics, lambdas, Swing, JDBC).; Bloch, \textit{Effective Java} (3e) (generics + lambdas).; Oracle Java Tutorials: Generics + Lambda Expressions + Swing + JDBC Basics (accessed \today).; Java SE API docs (e.g., \texttt{java.util.function}, \texttt{javax.swing}, \texttt{java.sql}) (accessed \today).}

\newcommand{\OOPSUnitVSources}{Schildt, \textit{Java: The Complete Reference} (12e) (Collections Framework + wrapper classes).; Oracle Java docs: Collections Framework overview (accessed \today).; Java SE API docs (\texttt{java.util} collections; wrapper classes in \texttt{java.lang}) (accessed \today).}

\newcommand{\OOPSUnitVISources}{Git SCM: \textit{Pro Git} (2e) (version control workflow), accessed \today.; GitHub Docs (repositories, issues, pull requests), accessed \today.; Oracle Java Tutorials: Swing + JDBC (accessed \today).; JUnit 5 User Guide (accessed \today).; Apache Maven Project: Getting Started (accessed \today).; Gradle User Manual: Getting Started (accessed \today).; UML class diagrams: UML 2.x overview/spec (accessed \today).}

\newcommand{\OOPSMiscWhyInterfacesSources}{Schildt, \textit{Java: The Complete Reference} (12e) (interfaces).; Bloch, \textit{Effective Java} (3e) (prefer interfaces where appropriate).; Oracle Java Tutorials: Interfaces (accessed \today).; Java SE API docs (e.g., \texttt{Comparable}, \texttt{Runnable}) (accessed \today).}



\title{Object Oriented Programming (CSEG1044)\\Unit 02 -- Lecture 02 Notes\\Constructors, `this`, Overloading, `static`, and GC Overview}
\author{Tofik Ali}
\date{\today}

\begin{document}
\maketitle
\tableofcontents

\section{Lecture Overview}
This lecture explains how objects are created and how to design classes cleanly:
\begin{itemize}
  \item constructors and constructor chaining,
  \item \texttt{this} vs \texttt{this(...)},
  \item method overloading,
  \item the \texttt{static} keyword,
  \item and a realistic overview of garbage collection (GC).
\end{itemize}
\notesources{\OOPSUnitIISources}

\section{Syllabus Alignment (Unit 02)}
\textbf{Unit 02:} constructors, \texttt{this}, overloading, \texttt{static}, garbage collection

\section{Learning Outcomes}
After this lecture, you should be able to:
\begin{itemize}
  \item Write default and parameterized constructors and use constructor chaining.
  \item Use \texttt{this} to refer to the current object and resolve shadowing.
  \item Overload methods correctly (and explain what is \emph{not} overloading).
  \item Use \texttt{static} fields/methods for class-level data/behavior.
  \item Explain what GC does and what it does \emph{not} guarantee.
\end{itemize}

\section{Key Terms (Quick)}
\begin{itemize}
  \item \textbf{Constructor}: special method to initialize an object (same name as class, no return type).
  \item \textbf{Shadowing}: parameter/local variable has same name as field (e.g., \texttt{name}).
  \item \textbf{\texttt{this}}: reference to the current object.
  \item \textbf{\texttt{this(...)} }: call another constructor in the same class (must be first statement).
  \item \textbf{Overloading}: same method name, different parameter list (compile-time selection).
  \item \textbf{\texttt{static}}: belongs to the class, shared across all objects.
  \item \textbf{Garbage collection}: automatic memory reclamation for unreachable objects.
\end{itemize}

\section{Detailed Notes}
\subsection{Constructors and constructor chaining}
Constructors are used to:
\begin{itemize}
  \item put an object into a valid initial state,
  \item enforce invariants early (validate inputs),
  \item and set default values when necessary.
\end{itemize}

\textbf{Default constructor rule:} If you do not write any constructor, Java provides a default no-arg constructor.
If you write \emph{any} constructor, Java does \emph{not} automatically create a default one.

\textbf{Constructor chaining:}
\begin{itemize}
  \item Use \texttt{this(...)} to call another constructor in the same class.
  \item \texttt{this(...)} must be the \textbf{first statement} in the constructor.
\end{itemize}

\subsection{\texttt{this} vs \texttt{this(...)}}
\begin{itemize}
  \item \texttt{this} refers to the current object (useful for field access and method calls).
  \item \texttt{this(...)} calls another constructor.
\end{itemize}

\textbf{Shadowing example:} if constructor parameter is \texttt{name}, the field is also \texttt{name}.
Use \texttt{this.name = name;} to distinguish.

\subsection{Method overloading (rules)}
Overloading means:
\begin{itemize}
  \item same method name,
  \item different parameter list (number/type/order).
\end{itemize}
Not overloading:
\begin{itemize}
  \item changing only the return type,
  \item changing only parameter names.
\end{itemize}

\paragraph{Why it matters.} Overloading lets you offer ``one concept'' in multiple forms (e.g., \texttt{print(int)} and \texttt{print(String)}).

\subsection{\texttt{static} keyword}
\textbf{Static field} is shared among all objects:
\begin{itemize}
  \item Use it for counters, constants, or shared configuration.
\end{itemize}

\textbf{Static method} belongs to the class:
\begin{itemize}
  \item You can call it as \texttt{ClassName.method()}.
  \item It cannot directly access instance fields (because there is no \texttt{this}).
\end{itemize}

\textbf{Utility class pattern:} when a class only has static methods (like \texttt{Math}), make the constructor \texttt{private} to prevent creating objects.

\subsection{Garbage collection overview}
GC frees memory of objects that are \textbf{no longer reachable} from any live reference.
Key points:
\begin{itemize}
  \item You cannot reliably force GC (\texttt{System.gc()} is only a suggestion).
  \item GC does not prevent all ``memory leaks'' if references are kept unnecessarily (e.g., in a static list).
  \item GC improves safety compared to manual memory management, but you still must close files, DB connections, etc.
\end{itemize}
\notesources{\OOPSUnitIISources}

\section{Worked Example: \texttt{Employee} (constructors + overloading + static)}
\begin{lstlisting}
public class Employee {
    private static int nextId = 1;

    private final int id;
    private final String name;
    private double salary;

    public Employee() {
        this("NA", 0.0); // constructor chaining
    }

    public Employee(String name, double salary) {
        if (name == null || name.isBlank()) throw new IllegalArgumentException("name");
        if (salary < 0) throw new IllegalArgumentException("salary");

        this.id = nextId++;
        this.name = name;
        this.salary = salary;
    }

    // Overloading: same name, different parameters
    public void raiseSalary(double percent) {
        raiseSalary(percent, 0.0);
    }

    public void raiseSalary(double percent, double bonus) {
        if (percent < 0) throw new IllegalArgumentException("percent");
        salary = salary + (salary * percent / 100.0) + bonus;
    }

    public static int getNextId() { return nextId; }
}
\end{lstlisting}

\section{In-Class Exercises (with brief solutions)}
\begin{enumerate}
  \item What is the difference between \texttt{this} and \texttt{this(...)}?\par
  \textbf{Answer:} \texttt{this} refers to current object; \texttt{this(...)} calls another constructor.

  \item Can we overload a method by changing only return type?\par
  \textbf{Answer:} No.

  \item Can a static method access an instance field directly?\par
  \textbf{Answer:} No, because there is no \texttt{this}. It needs an object reference.
\end{enumerate}

\section{Self-Study Practice (no solutions)}
\begin{enumerate}
  \item Create a \texttt{Point} class with 3 constructors (0-arg, 2-arg, copy constructor) using chaining.
  \item Overload a method \texttt{area(...)} for a square (1 param) and rectangle (2 params).
  \item Write a class-level counter using a static field and print how many objects were created.
\end{enumerate}

\section{Checkpoint Questions (with sample answers)}
\textbf{Q1.} Why do we use constructor chaining?\par
\textbf{Sample answer:} It reduces duplicate initialization code. One constructor can call the main constructor that performs validation and sets fields consistently.

\vspace{0.6em}
\textbf{Q2.} What does garbage collection do?\par
\textbf{Sample answer:} GC reclaims memory of objects that are no longer reachable. It does not guarantee immediate collection and does not automatically close external resources like files or DB connections.

\end{document}
